% Generated by Listas/make_statements.py from pratica-lecture-04.md, item 11. Do not edit.
\textbf{Source.} Tadelis (2013), \emph{Game Theory: An Introduction}, Exercise 8.12 ``Agenda Setting'', p.~173 {[}PDF 192{]}. Reproduced verbatim. Answer in the source: none in the book (the preface, p.~xiii, says about half are solved on Princeton UP's website; not checked).



\begin{quote}

\textbf{Agenda Setting:} An agenda-setting game is described as follows. The ``issue space'' (set of possible policies) is an interval \(X = [0, 5]\). An agenda setter (player 1) proposes an alternative \(x \in X\) against the status quo \(q = 4\). After player 1 proposes \(x\), the legislator (player 2) observes the proposal and selects between the proposal \(x\) and the status quo \(q\). Player 1's most-preferred policy is 1, and for any final policy \(y \in X\) his payoff is given by



\[v_1(y) = 10 - |y - 1|,\]



where \(|y - 1|\) denotes the absolute value of \((y - 1)\). Player 2's most preferred policy is 3, and for any final policy \(y \in X\) his payoff is given by



\[v_2(y) = 10 - |y - 3|.\]



That is, each player prefers policies that are closer to his most-preferred policy.



\begin{enumerate}

\def\labelenumi{\alph{enumi}.}

\item

  Write the game down as a normal-form game. Is this a game of perfect or imperfect information?

\item

  Find a subgame-perfect equilibrium for this game. Is it unique?

\item

  Find a Nash equilibrium that is not subgame perfect. Is it unique? If yes, explain. If not, show all the Nash equilibria for this game.

\end{enumerate}

\end{quote}

